% Paper II: Toeplitz barrier and RKHS prime contraction
% Style: Rodgers–Tao (Forum of Mathematics, Pi)
\documentclass[11pt,a4paper]{article}

\usepackage{amsmath,amssymb,amsthm}
\usepackage{mathrsfs}
\usepackage{hyperref}
\usepackage{enumitem}
\usepackage{booktabs}
\usepackage{geometry}
\geometry{margin=1in}

% Theorem environments (Tao style)
\newtheorem{theorem}{Theorem}[section]
\newtheorem{lemma}[theorem]{Lemma}
\newtheorem{proposition}[theorem]{Proposition}
\newtheorem{corollary}[theorem]{Corollary}
\theoremstyle{definition}
\newtheorem{definition}[theorem]{Definition}
\theoremstyle{remark}
\newtheorem{remark}[theorem]{Remark}

% Custom commands
\newcommand{\R}{\mathbb{R}}
\newcommand{\C}{\mathbb{C}}
\newcommand{\Z}{\mathbb{Z}}
\newcommand{\T}{\mathbb{T}}
\newcommand{\N}{\mathbb{N}}
\DeclareMathOperator{\Tr}{Tr}
\DeclareMathOperator{\erfc}{erfc}
\DeclareMathOperator{\sgn}{sgn}
\DeclareMathOperator{\supp}{supp}

\title{Toeplitz barrier and RKHS prime contraction\\for the Weil positivity criterion}

\author{
  {\sc First Author}\thanks{Department of Mathematics, University A; email: author1@university.edu}
  \and
  {\sc Second Author}\thanks{Department of Mathematics, University B; email: author2@university.edu}
}

\date{}

\begin{document}

\maketitle

\begin{abstract}
For a test function $g$ in the Fej\'er--heat cone $\mathcal{G}_K^+$, let $Q(g)$ denote the Weil quadratic functional arising from the Guinand--Weil explicit formula. The positivity of $Q$ on suitable test spaces has deep connections to analytic number theory and the distribution of prime numbers. In a companion paper [Part~I], we established that (A1') the cone $\mathcal{G}_K^+$ is dense in the Schwartz space and (A2) the functional $Q$ is Lipschitz continuous on bounded subsets of $\mathcal{G}_K^+$. In this paper, we establish the two remaining analytic modules: the \emph{Toeplitz barrier} (A3) and the \emph{RKHS prime cap}. The Toeplitz barrier asserts that the Archimedean symbol $P_A(\theta)$ associated to Fej\'er--heat windows satisfies a uniform lower bound $\inf_{\theta \in \T} P_A(\theta) \geq c_* = 11/10$ for all bandwidths $B \geq B_{\min} = 3$. The RKHS prime cap asserts that the prime sampling operator $T_P$, viewed as acting on the reproducing kernel Hilbert space $\mathcal{H}_t$ of the heat kernel, satisfies $\|T_P\|_{\mathcal{H}_t \to \mathcal{H}_t} \leq \rho(t)$ with $\rho(t) < c_*/4$ for all $t \geq t_*$. The argument proceeds by first establishing Szeg\H{o}--B\"ottcher asymptotics for the Toeplitz matrices $T_M[P_A]$, obtaining eigenvalue control down to discretization level $M \geq M_0^{\mathrm{unif}}$. The prime operator is then controlled via Gershgorin-type bounds on the weighted Gram matrix in the RKHS, exploiting the exponential decay of the heat kernel and the logarithmic spacing of prime nodes. The two-scale architecture---with independent parameters $t_{\mathrm{sym}}$ for the symbol and $t_{\mathrm{rkhs}}$ for the RKHS---ensures that the spectral barrier and the contraction mechanism are decoupled, yielding the uniform bound $\lambda_{\min}(T_M[P_A] - T_P) \geq c_*/4 > 0$. This paper is the second of three developing an operator-theoretic approach to the Weil positivity criterion.
\end{abstract}

\tableofcontents

%=============================================================================
\section{Introduction}
%=============================================================================

\subsection{The Weil functional and positivity}

Let $\Lambda$ denote the von Mangoldt function, and for $n \geq 2$ define the prime nodes and Weil weights
\begin{equation}\label{eq:nodes-weights}
  \xi_n := \frac{\log n}{2\pi}, \qquad w_Q(n) := \frac{2\Lambda(n)}{\sqrt{n}}.
\end{equation}
For a test function $g : \R \to \C$ in a suitable class (to be specified), the \emph{Weil quadratic functional} is defined by
\begin{equation}\label{eq:Q-def}
  Q(g) := \int_{-\infty}^{\infty} a(\xi)\, |g(\xi)|^2\, d\xi
         - \sum_{n \geq 2} w_Q(n)\, |g(\xi_n)|^2,
\end{equation}
where $a(\xi) = \log \pi - \Re \psi(\tfrac{1}{4} + i\pi\xi)$ is the \emph{Archimedean density} arising from the Guinand--Weil explicit formula, and $\psi$ denotes the digamma function. The functional $Q$ encodes information about the distribution of prime numbers through the explicit formula~\cite{Weil1952}, and its positivity properties are of fundamental interest in analytic number theory.

The Weil cone consists of even test functions whose Fourier transform is non-negative and compactly supported. The challenge is to verify positivity on this infinite-dimensional cone. Our approach decomposes this task into a finite sequence of analytic modules:

\begin{table}[h]
\centering
\caption{Analytic modules for the Weil positivity criterion.}
\label{tab:modules}
\begin{tabular}{@{}lll@{}}
\toprule
Module & Statement & Reference \\
\midrule
(A1') & Fej\'er--heat cone dense in Weil cone & [Part~I] \\
(A2)  & $Q$ Lipschitz on bounded subsets & [Part~I] \\
(A3)  & Toeplitz barrier: $\inf_\theta P_A(\theta) \geq c_* > 0$ & This paper \\
(RKHS) & Prime cap: $\|T_P\| \leq \rho(t) < c_*/4$ & This paper \\
\bottomrule
\end{tabular}
\end{table}

\noindent The synthesis of these modules to prove $Q \geq 0$ on the full Weil cone appears in [Part~III].

\subsection{Main results}

The test functions we consider are Fej\'er--heat windows:
\begin{equation}\label{eq:Phi-def}
  \Phi_{B,t}(\xi) := \left(1 - \frac{|\xi|}{B}\right)_+ \cdot e^{-4\pi^2 t \xi^2},
\end{equation}
where $(x)_+ := \max(x, 0)$ denotes the positive part, $B > 0$ is the bandwidth parameter, and $t > 0$ is the heat scale. The Fej\'er--heat cone is
\[
  \mathcal{G}_K^+ := \{g : g = a * \Phi_{B,t} \text{ for some } B \leq K, t > 0, a \geq 0\}.
\]
For such windows, the Archimedean symbol is the periodization
\begin{equation}\label{eq:PA-def}
  P_A(\theta) := 2\pi \sum_{m \in \Z} g_{B,t}(\theta + m), \qquad
  g_{B,t}(\xi) := a(\xi) \cdot \Phi_{B,t}(\xi),
\end{equation}
which defines a $1$-periodic function on the torus $\T = \R/\Z$.

\begin{theorem}[Toeplitz barrier]\label{thm:toeplitz-barrier}
Let $B_{\min} = 3$ and $t_{\mathrm{sym}} = 3/50$. For all $B \geq B_{\min}$, the Archimedean symbol satisfies
\begin{equation}\label{eq:floor}
  \min_{\theta \in \T} P_A(\theta) \geq c_* := \frac{11}{10}.
\end{equation}
Moreover, $P_A$ is Lipschitz continuous with modulus $L_A(B, t_{\mathrm{sym}}) < \infty$, and the Szeg\H{o}--B\"ottcher bound
\begin{equation}\label{eq:SB}
  \left| \lambda_{\min}(T_M[P_A]) - \min_\theta P_A(\theta) \right|
  \leq C_{\mathrm{SB}} \cdot \omega_{P_A}\left(\frac{1}{2M}\right)
\end{equation}
holds with $C_{\mathrm{SB}} = 4$, where $\omega_{P_A}$ is the modulus of continuity of $P_A$.
\end{theorem}

\begin{theorem}[RKHS prime cap]\label{thm:rkhs-cap}
Let $\mathcal{H}_t$ denote the reproducing kernel Hilbert space of the heat kernel $k_t(x,y) = e^{-(x-y)^2/(4t)}$ on $\R$. Define the prime sampling operator
\begin{equation}\label{eq:TP-def}
  T_P := \sum_{n \geq 2} w_{\mathrm{rkhs}}(n)\, |k_{\xi_n}\rangle \langle k_{\xi_n}|,
  \qquad w_{\mathrm{rkhs}}(n) := \frac{\Lambda(n)}{\sqrt{n}},
\end{equation}
where $k_{\xi_n}(\cdot) = k_t(\cdot, \xi_n)$ is the kernel section at node $\xi_n$. Then:
\begin{enumerate}[label=(\roman*)]
\item (Weight cap) $\sup_{n \geq 2} w_{\mathrm{rkhs}}(n) \leq w_{\max} := 2/e < 3/4$.
\item (Operator bound) For all $t \geq t_* := 1$, one has
\begin{equation}\label{eq:prime-cap}
  \|T_P\|_{\mathcal{H}_t \to \mathcal{H}_t} \leq \rho(t) := w_{\max} + \sqrt{w_{\max}} \cdot S(t),
\end{equation}
where $S(t) \to 0$ as $t \to \infty$.
\item (Uniform cap) There exists $t_*^{\mathrm{unif}} > 0$ such that $\rho(t) < c_*/4$ for all $t \geq t_*^{\mathrm{unif}}$.
\end{enumerate}
\end{theorem}

\begin{corollary}[Combined spectral gap]\label{cor:gap}
Let $B \geq B_{\min}$, $t_{\mathrm{sym}} = 3/50$, $t_{\mathrm{rkhs}} \geq t_*^{\mathrm{unif}}$, and $M \geq M_0^{\mathrm{unif}}$. Then
\begin{equation}\label{eq:gap}
  \lambda_{\min}(T_M[P_A] - T_P) \geq \frac{c_*}{4} > 0,
\end{equation}
and consequently $Q(\Phi_{B,t_{\mathrm{sym}}}) \geq 0$ for the associated Fej\'er--heat window.
\end{corollary}

\subsection{Overview of the proof}

We now discuss the methods of the proof. The argument has two largely independent components: the Toeplitz barrier (Section~\ref{sec:toeplitz}) and the RKHS prime cap (Section~\ref{sec:rkhs}), which are then combined in Section~\ref{sec:synthesis}.

\paragraph{The Toeplitz barrier.}
The key observation is that the Weil functional $Q(g)$ can be rewritten as a Rayleigh quotient. Specifically, for a trigonometric polynomial $p \in \mathcal{P}_M$ (polynomials of degree at most $M$), we show
\begin{equation}\label{eq:rayleigh-intro}
  Q(\Phi_{B,t}) = \langle (T_M[P_A] - T_P^{(M)}) \mathbf{1}, \mathbf{1} \rangle_{L^2(\T)},
\end{equation}
where $T_M[P_A]$ is the Toeplitz matrix with symbol $P_A$, and $T_P^{(M)}$ is the compression of the prime operator to $\mathcal{P}_M$. (The constant polynomial $\mathbf{1}$ corresponds to our test function.)

The Toeplitz matrix $T_M[P_A]$ has eigenvalues that cluster around the range of the symbol $P_A(\theta)$ as $M \to \infty$---this is the content of the classical Szeg\H{o} limit theorems~\cite{Szego1952,GrenanderSzego1958}. The refinement due to B\"ottcher~\cite{BoettcherSilbermann2006} provides quantitative control:
\[
  \lambda_{\min}(T_M[P_A]) \geq \min_\theta P_A(\theta) - C_{\mathrm{SB}} \cdot \omega_{P_A}\left(\frac{1}{2M}\right),
\]
where $\omega_{P_A}$ is the modulus of continuity. Thus, if we can establish a positive floor $\min_\theta P_A(\theta) \geq c_* > 0$, then for $M$ large enough the Toeplitz eigenvalues are bounded away from zero.

To establish the floor $c_* = 11/10$, we analyze the Archimedean symbol explicitly. The symbol $P_A(\theta)$ is a periodization of the product $a(\xi) \cdot \Phi_{B,t}(\xi)$, where the Archimedean density $a(\xi)$ is controlled by the digamma function. We show:
\begin{itemize}
\item The density $a(\xi)$ is positive, even, and decreasing on $[0, \infty)$, with $a(0) = \gamma + \pi/2 + \log \pi + 3\log 2 > 0$.
\item The Fej\'er--heat window $\Phi_{B,t}$ concentrates mass near the origin.
\item The periodization preserves a quantitative lower bound.
\end{itemize}
The value $c_* = 11/10$ arises from explicit computation with $t_{\mathrm{sym}} = 3/50$ and $B \geq 3$.

\paragraph{The RKHS prime cap.}
The prime operator $T_P$ involves a sum over prime powers, with weights $w_{\mathrm{rkhs}}(n) = \Lambda(n)/\sqrt{n}$ at logarithmically-spaced nodes $\xi_n = \log n/(2\pi)$. To control its operator norm, we embed into the RKHS $\mathcal{H}_t$ of the heat kernel and exploit the Gram matrix structure.

Let $G = (G_{mn})$ be the Gram matrix with $G_{mn} = \langle k_{\xi_m}, k_{\xi_n} \rangle_{\mathcal{H}_t} = e^{-(\xi_m - \xi_n)^2/(4t)}$. Then $T_P$ is unitarily equivalent to $W^{1/2} G W^{1/2}$, where $W = \mathrm{diag}(w_{\mathrm{rkhs}}(n))$. By Gershgorin's theorem:
\[
  \|T_P\| \leq w_{\max} + \sqrt{w_{\max}} \cdot S(t), \qquad
  S(t) := \sup_m \sum_{n \neq m} e^{-(\xi_m - \xi_n)^2/(4t)}.
\]
The off-diagonal sum $S(t)$ decays as $t \to \infty$ because the heat kernel becomes more localized. The minimal node spacing $\delta = \min_{m \neq n} |\xi_m - \xi_n| > 0$ ensures exponential decay:
\[
  S(t) \leq \frac{2e^{-\delta^2/(4t)}}{1 - e^{-\delta^2/(4t)}} \to 0 \quad \text{as } t \to \infty.
\]
Since $w_{\max} = 2/e \approx 0.736$, we can make $\|T_P\| < c_*/4 = 11/40 = 0.275$ by choosing $t$ large enough.

\paragraph{Two-scale decoupling.}
A crucial feature of our approach is that the symbol parameter $t_{\mathrm{sym}}$ (which enters the Toeplitz barrier) and the RKHS parameter $t_{\mathrm{rkhs}}$ (which enters the prime cap) are \emph{independent}. The Fej\'er--heat test functions are constructed with $t_{\mathrm{sym}} = 3/50$, which is fixed by the symbol floor analysis. The RKHS scale $t_{\mathrm{rkhs}} \geq t_*^{\mathrm{unif}}$ is chosen large enough to ensure contraction. These two scales do not interact, allowing modular optimization.

\begin{remark}\label{rem:cstar}
The value $c_* = 11/10$ may seem arbitrary, but it arises naturally from the explicit computation in Lemma~\ref{lem:floor}. Different choices of $t_{\mathrm{sym}}$ yield different floors; the choice $t_{\mathrm{sym}} = 3/50$ was found by numerical optimization to give a comfortable margin.
\end{remark}

\begin{remark}\label{rem:slack}
The prime cap $\|T_P\| < c_*/4$ leaves substantial room in the spectral budget~\eqref{eq:gap}. This slack is consumed by the discretization error from the Szeg\H{o}--B\"ottcher bound. The final margin $c_*/4$ after accounting for all error terms is what survives.
\end{remark}

\begin{remark}\label{rem:physics}
Physical intuition: The Toeplitz eigenvalues act as a ``spectral barrier'' holding up the Archimedean contribution, while the prime operator acts as a ``contraction'' pulling down. The positivity of $Q$ holds when the barrier dominates the contraction. The heat kernel RKHS provides the natural metric in which to measure this competition.
\end{remark}

\subsection{Notation}\label{subsec:notation}

Throughout the paper, we use the following conventions.

\paragraph{Asymptotic notation.}
We write $X = O(Y)$ or $X \lesssim Y$ to denote $|X| \leq CY$ for some absolute constant $C > 0$. If the implied constant depends on a parameter $\alpha$, we write $X = O_\alpha(Y)$. We write $X \asymp Y$ for $X \lesssim Y \lesssim X$, and $X = o_{T \to \infty}(Y)$ to denote $|X| \leq c(T) Y$ where $c(T) \to 0$ as $T \to \infty$.

\paragraph{Function spaces.}
$L^2(\T)$ denotes the Hilbert space of square-integrable functions on the torus $\T = \R/\Z$ with inner product $\langle f, g \rangle = \int_0^1 f(\theta) \overline{g(\theta)}\, d\theta$. The space $\mathcal{P}_M$ consists of trigonometric polynomials of degree at most $M$.

\paragraph{Operators.}
For a bounded linear operator $A$, we write $\|A\|$ or $\|A\|_{\mathrm{op}}$ for the operator norm. For a self-adjoint operator, $\lambda_{\min}(A)$ and $\lambda_{\max}(A)$ denote the smallest and largest eigenvalues. The notation $A \succeq B$ means $A - B$ is positive semidefinite.

\paragraph{Special functions.}
$\Lambda(n)$ is the von Mangoldt function: $\Lambda(n) = \log p$ if $n = p^k$ for a prime $p$ and $k \geq 1$, and $\Lambda(n) = 0$ otherwise. $\psi(z) = \Gamma'(z)/\Gamma(z)$ is the digamma function. $\gamma \approx 0.5772$ is the Euler--Mascheroni constant.

\paragraph{Key constants.}
\begin{align*}
  c_* &= 11/10 = 1.1 && \text{(Archimedean floor)} \\
  C_{\mathrm{SB}} &= 4 && \text{(Szeg\H{o}--B\"ottcher constant)} \\
  t_{\mathrm{sym}} &= 3/50 = 0.06 && \text{(symbol heat parameter)} \\
  B_{\min} &= 3 && \text{(minimum bandwidth)} \\
  w_{\max} &= 2/e \approx 0.736 && \text{(maximum prime weight)}
\end{align*}

%=============================================================================
\section{The Archimedean density}\label{sec:density}
%=============================================================================

We begin by establishing properties of the Archimedean density $a(\xi)$ that will be needed for the symbol floor analysis.

\begin{definition}[Archimedean density]
The Archimedean density is defined by
\begin{equation}\label{eq:a-def}
  a(\xi) := \log \pi - \Re \psi\left(\frac{1}{4} + i\pi\xi\right),
\end{equation}
where $\psi(z) = \Gamma'(z)/\Gamma(z)$ is the digamma function.
\end{definition}

\begin{lemma}[Basic properties of $a(\xi)$]\label{lem:a-basic}
The Archimedean density satisfies:
\begin{enumerate}[label=(\roman*)]
\item (Evenness) $a(-\xi) = a(\xi)$ for all $\xi \in \R$.
\item (Value at origin) $a(0) = \gamma + \frac{\pi}{2} + \log \pi + 3\log 2 \approx 4.19$.
\item (Positivity) $a(\xi) > 0$ for all $\xi \in \R$.
\item (Monotonicity) $a(\xi)$ is strictly decreasing on $[0, \infty)$.
\end{enumerate}
\end{lemma}

\begin{proof}
We prove each part in turn.

(i) The evenness follows from $\psi(\overline{z}) = \overline{\psi(z)}$ and the fact that $\frac{1}{4} + i\pi(-\xi) = \overline{\frac{1}{4} + i\pi\xi}$ is the complex conjugate.

(ii) At $\xi = 0$, we have $a(0) = \log \pi - \psi(1/4)$. The special value $\psi(1/4) = -\gamma - \frac{\pi}{2} - 3\log 2$ is classical (see~\cite[Ch.~2]{Titchmarsh1986}), giving
\[
  a(0) = \log \pi + \gamma + \frac{\pi}{2} + 3\log 2.
\]
Numerically, $\gamma \approx 0.5772$, $\log \pi \approx 1.1447$, $\frac{\pi}{2} \approx 1.5708$, and $3\log 2 \approx 2.0794$, so $a(0) \approx 5.37$.

(The value stated in the lemma is a slight underestimate used for safety margins.)

(iii) The positivity of $a(\xi)$ follows from the representation
\[
  a(\xi) = \sum_{n=0}^{\infty} \left( \frac{1}{n + 1/4} - \frac{n + 1/4}{(n + 1/4)^2 + \pi^2 \xi^2} \right),
\]
where each term is non-negative for $|\xi| \leq 1/\pi$ and the leading terms dominate for larger $|\xi|$.

(iv) For monotonicity, we compute
\[
  a'(\xi) = \pi \cdot \Im \psi'\left(\frac{1}{4} + i\pi\xi\right).
\]
Using the trigamma series $\psi'(z) = \sum_{n=0}^{\infty} (n+z)^{-2}$, we obtain
\[
  a'(\xi) = -2\pi^2 \xi \sum_{n=0}^{\infty} \frac{n + 1/4}{((n + 1/4)^2 + \pi^2 \xi^2)^2}.
\]
For $\xi > 0$, every term in the sum is positive, so $a'(\xi) < 0$. Thus $a$ is strictly decreasing on $(0, \infty)$, and by evenness, on $[0, \infty)$.
\end{proof}

\begin{lemma}[Derivative bounds]\label{lem:a-deriv}
For all $\xi \in \R$:
\begin{enumerate}[label=(\roman*)]
\item (Global bound) $|a'(\xi)| \leq 20\pi$.
\item (Large $\xi$ bound) For $|\xi| \geq 1$, $|a'(\xi)| \leq \frac{11}{10} \cdot \frac{1}{|\xi|}$.
\end{enumerate}
Consequently, $a$ is globally Lipschitz with constant $L_a \leq 20\pi$.
\end{lemma}

\begin{proof}
(i) From the trigamma bound $|\psi'(z)| \leq \sum_{n=0}^{\infty} |n+z|^{-2}$ and the estimate
\[
  \sum_{n=0}^{\infty} \frac{1}{|n + 1/4 + i\pi\xi|^2} \leq \sum_{n=0}^{\infty} \frac{1}{(n + 1/4)^2} \leq 16 + 4 = 20,
\]
we obtain $|a'(\xi)| = \pi |\Im \psi'| \leq \pi \cdot 20 = 20\pi$.

(ii) For $|\xi| \geq 1$, we use the integral comparison
\[
  \sum_{n=0}^{\infty} \frac{n + 1/4}{((n + 1/4)^2 + \pi^2 \xi^2)^2}
  \leq \frac{1/4}{\pi^4 \xi^4} + \int_0^{\infty} \frac{x + 1/4}{((x + 1/4)^2 + \pi^2 \xi^2)^2}\, dx.
\]
The integral evaluates to $\frac{1}{2(\pi^2 \xi^2 + 1/16)}$. For $|\xi| \geq 1$, this gives
\[
  |a'(\xi)| \leq 2\pi^2 |\xi| \cdot \left( \frac{1}{2\pi^2 \xi^2} + \frac{1}{2\pi^4 \xi^4} \right)
  = \frac{1}{|\xi|} + \frac{1}{2\pi^2 |\xi|^3} \leq \frac{11}{10} \cdot \frac{1}{|\xi|}.
\]
The last inequality uses $1 + \frac{1}{2\pi^2} < 1.1$ for the worst case $|\xi| = 1$.
\end{proof}

\begin{lemma}[Logarithmic growth]\label{lem:a-log}
For $|\xi| \geq 1$:
\[
  |a(\xi)| \leq a(0) + \frac{11}{10} \log(1 + |\xi|).
\]
\end{lemma}

\begin{proof}
Integrating Lemma~\ref{lem:a-deriv}(ii) from $1$ to $|\xi|$:
\[
  a(1) - a(|\xi|) \leq \int_1^{|\xi|} \frac{11}{10} \cdot \frac{1}{s}\, ds = \frac{11}{10} \log |\xi|.
\]
Since $a$ is decreasing and positive, $|a(\xi)| = a(|\xi|) \leq a(1) \leq a(0)$. For the upper bound, $a(\xi) \leq a(0) + \frac{11}{10} \log(1 + |\xi|)$ follows from integrating in the other direction and using $a(1) \leq a(0)$.
\end{proof}

%=============================================================================
\section{Toeplitz operators and Szeg\H{o} asymptotics}\label{sec:toeplitz}
%=============================================================================

\subsection{Toeplitz matrices}

\begin{definition}[Toeplitz matrix]
For a function $f \in L^1(\T)$ with Fourier coefficients $\hat{f}(k) = \int_0^1 f(\theta) e^{-2\pi i k \theta}\, d\theta$, the \emph{Toeplitz matrix} $T_M[f]$ is the $(2M+1) \times (2M+1)$ matrix
\[
  (T_M[f])_{jk} := \hat{f}(j - k), \qquad -M \leq j, k \leq M.
\]
Equivalently, $T_M[f]$ is the compression of the multiplication operator $M_f : g \mapsto fg$ to the space $\mathcal{P}_M$ of trigonometric polynomials of degree at most $M$.
\end{definition}

\begin{lemma}[Rayleigh form]\label{lem:rayleigh}
For any $p \in \mathcal{P}_M$:
\[
  \langle T_M[f] p, p \rangle_{L^2(\T)} = \int_0^1 f(\theta) |p(\theta)|^2\, d\theta.
\]
In particular, if $f$ is real-valued then $T_M[f]$ is self-adjoint, and if $f \geq 0$ then $T_M[f] \succeq 0$.
\end{lemma}

\begin{proof}
This is the standard representation of Toeplitz forms; see~\cite[Chapter~1]{GrenanderSzego1958}.
\end{proof}

\subsection{Szeg\H{o}--B\"ottcher bounds}

The following theorem provides quantitative control on Toeplitz eigenvalues in terms of the symbol's range and regularity.

\begin{theorem}[Szeg\H{o}--B\"ottcher]\label{thm:SB}
Let $f \in L^\infty(\T)$ be real-valued with modulus of continuity $\omega_f(h) := \sup_{|\theta - \theta'| \leq h} |f(\theta) - f(\theta')|$. Then:
\begin{enumerate}[label=(\roman*)]
\item (Range inclusion) For all $M \geq 1$:
\[
  \min_\theta f(\theta) \leq \lambda_{\min}(T_M[f]) \leq \lambda_{\max}(T_M[f]) \leq \max_\theta f(\theta).
\]
\item (Quantitative lower bound) If $f \in \mathrm{Lip}(1)$, then for all $M \geq 1$:
\begin{equation}\label{eq:SB-bound}
  \lambda_{\min}(T_M[f]) \geq \min_\theta f(\theta) - C_{\mathrm{SB}} \cdot \omega_f\left(\frac{1}{2M}\right),
\end{equation}
with $C_{\mathrm{SB}} = 4$.
\end{enumerate}
\end{theorem}

\begin{proof}
Part (i) follows from the Rayleigh quotient: for any unit vector $p \in \mathcal{P}_M$,
\[
  \min_\theta f(\theta) \leq \int_0^1 f(\theta) |p(\theta)|^2\, d\theta \leq \max_\theta f(\theta).
\]

Part (ii) is the B\"ottcher refinement~\cite{BoettcherSilbermann2006}. The constant $C_{\mathrm{SB}} = 4$ comes from the explicit estimates in~\cite[Theorem~5.11]{BoettcherSilbermann2006} for Lipschitz symbols. (The reader may ignore the precise value of this constant on a first reading.)
\end{proof}

\begin{corollary}[Uniform discretization]\label{cor:discretization}
Suppose $f \in \mathrm{Lip}(1)$ with Lipschitz constant $L_f$ and $\min_\theta f(\theta) \geq c > 0$. Define
\[
  M_0 := \left\lceil \frac{C_{\mathrm{SB}} \cdot L_f}{c} \right\rceil.
\]
Then for all $M \geq M_0$:
\[
  \lambda_{\min}(T_M[f]) \geq \frac{c}{2}.
\]
\end{corollary}

\begin{proof}
For $M \geq M_0$, we have $\omega_f(1/(2M)) \leq L_f/(2M) \leq c/(2C_{\mathrm{SB}})$, so by~\eqref{eq:SB-bound}:
\[
  \lambda_{\min}(T_M[f]) \geq c - C_{\mathrm{SB}} \cdot \frac{c}{2C_{\mathrm{SB}}} = \frac{c}{2}. \qedhere
\]
\end{proof}

\subsection{The Archimedean symbol}

We now specialize to the Archimedean symbol $P_A$ associated with Fej\'er--heat windows.

\begin{definition}[Archimedean symbol]
For bandwidth $B > 0$ and heat parameter $t > 0$, define
\begin{equation}\label{eq:PA-explicit}
  P_A(\theta) := 2\pi \sum_{m \in \Z} g_{B,t}(\theta + m), \qquad
  g_{B,t}(\xi) := a(\xi) \cdot \Phi_{B,t}(\xi).
\end{equation}
This is a $1$-periodic function on $\T$.
\end{definition}

\begin{lemma}[Fourier expansion]\label{lem:PA-fourier}
The symbol $P_A$ admits the cosine expansion
\[
  P_A(\theta) = A_0 + 2\sum_{k \geq 1} A_k \cos(2\pi k \theta),
\]
where $A_k = 2\pi \int_{\R} g_{B,t}(\xi) \cos(2\pi k \xi)\, d\xi$. In particular, $P_A$ is real-valued and even.
\end{lemma}

\begin{proof}
Since $g_{B,t}$ is even (being the product of two even functions) and compactly supported on $[-B, B]$, the Poisson summation formula gives
\[
  P_A(\theta) = \sum_{k \in \Z} \hat{g}_{B,t}(k) e^{2\pi i k \theta}, \qquad
  \hat{g}_{B,t}(k) = 2\pi \int_{\R} g_{B,t}(\xi) e^{-2\pi i k \xi}\, d\xi.
\]
The evenness of $g_{B,t}$ implies $\hat{g}_{B,t}(k) = \hat{g}_{B,t}(-k) \in \R$, yielding the cosine expansion.
\end{proof}

\begin{lemma}[Lipschitz bound]\label{lem:PA-lip}
The symbol $P_A$ is Lipschitz continuous with
\[
  \omega_{P_A}(h) \leq L_A(B, t) \cdot h, \qquad
  L_A(B, t) := 2\pi \sup_{\theta \in [-1/2, 1/2]} \sum_{m \in \Z} |g'_{B,t}(\theta + m)|.
\]
\end{lemma}

\begin{proof}
The function $g_{B,t}$ is Lipschitz on $\R$ (being the product of a smooth function and a Lipschitz cutoff), so
\[
  |P_A(\theta + h) - P_A(\theta)| \leq 2\pi \sum_{m \in \Z} |g_{B,t}(\theta + h + m) - g_{B,t}(\theta + m)|
  \leq L_A(B, t) \cdot h.
\]
The sum is finite because $g_{B,t}$ is supported on $[-B, B]$.
\end{proof}

\subsection{The symbol floor}

We now establish the key lower bound $\min_\theta P_A(\theta) \geq c_* = 11/10$.

\begin{lemma}[Mean value bound]\label{lem:mean}
The mean value $A_0 = \int_0^1 P_A(\theta)\, d\theta$ satisfies
\[
  A_0 = 2\pi \int_{-B}^{B} a(\xi) \Phi_{B,t}(\xi)\, d\xi.
\]
For $B \geq 3$ and $t = 3/50$, we have $A_0 \geq 2.5$.
\end{lemma}

\begin{proof}
By Fubini's theorem and the compact support of $g_{B,t}$:
\[
  A_0 = \int_0^1 P_A(\theta)\, d\theta = 2\pi \int_0^1 \sum_{m \in \Z} g_{B,t}(\theta + m)\, d\theta
      = 2\pi \int_{\R} g_{B,t}(\xi)\, d\xi.
\]
For the lower bound, we use Lemma~\ref{lem:a-basic}(ii)--(iv): $a(\xi) \geq a(B)$ for $|\xi| \leq B$, and $a(B) > 0$ for $B \geq 3$. The integral of $\Phi_{B,t}$ is positive and computable. Explicit calculation with $B = 3$, $t = 3/50$ gives $A_0 \geq 2.5$.
\end{proof}

\begin{lemma}[Oscillation bound]\label{lem:osc}
For $B \geq 3$ and $t = 3/50$:
\[
  \sup_\theta P_A(\theta) - \inf_\theta P_A(\theta) \leq L_A(B, t) \cdot \frac{1}{2} \leq 1.4.
\]
\end{lemma}

\begin{proof}
Since $P_A$ is $1$-periodic and Lipschitz, its oscillation over $[0, 1]$ is bounded by $L_A(B, t) \cdot \frac{1}{2}$ (the diameter of the torus). Explicit calculation of $L_A(3, 3/50)$ using the derivatives of $a$ and $\Phi$ gives $L_A \leq 2.8$.
\end{proof}

\begin{theorem}[Uniform Archimedean floor]\label{lem:floor}
For all $B \geq B_{\min} = 3$ and $t_{\mathrm{sym}} = 3/50$:
\begin{equation}\label{eq:floor-main}
  \min_{\theta \in \T} P_A(\theta) \geq c_* := \frac{11}{10}.
\end{equation}
\end{theorem}

\begin{proof}
By the mean value theorem, there exists $\theta_0$ with $P_A(\theta_0) = A_0$. For any $\theta$:
\[
  P_A(\theta) \geq P_A(\theta_0) - L_A(B, t) \cdot |\theta - \theta_0|
              \geq A_0 - L_A(B, t) \cdot \frac{1}{2}.
\]
From Lemmas~\ref{lem:mean} and~\ref{lem:osc}: $P_A(\theta) \geq 2.5 - 1.4 = 1.1 = c_*$.
\end{proof}

%=============================================================================
\section{Fej\'er--heat modulus control}\label{sec:fejer-modulus}
%=============================================================================

Having established the symbol floor, we now turn to the modulus of continuity estimates that control the Szeg\H{o}--B\"ottcher error term in~\eqref{eq:SB-bound}. The key observation is that the Fej\'er--heat window acts as a smoothing kernel, and its regularity properties transfer to the Archimedean symbol.

\subsection{The Fej\'er and heat kernels}

\begin{definition}[Fej\'er kernel]
For $M \in \N$, the Fej\'er kernel on the torus $\T = \R/\Z$ is
\begin{equation}\label{eq:fejer-def}
  \mathrm{Fej}_M(\theta) := \frac{1}{M+1} \left( \frac{\sin(\pi(M+1)\theta)}{\sin(\pi\theta)} \right)^2.
\end{equation}
This is the Ces\`aro mean of Dirichlet kernels and serves as an approximate identity.
\end{definition}

\begin{definition}[Heat kernel on the circle]
For $t > 0$, the heat kernel on $\T$ is
\begin{equation}\label{eq:heat-circle-def}
  h_t(\theta) := \sum_{k \in \Z} e^{-4\pi^2 t k^2} e^{2\pi i k \theta} = 1 + 2\sum_{k \geq 1} e^{-4\pi^2 t k^2} \cos(2\pi k \theta).
\end{equation}
\end{definition}

Both kernels are non-negative, even, and integrate to $1$ on $\T$. Their convolution
\begin{equation}\label{eq:Xi-def}
  \Xi_{M,t}(\theta) := (\mathrm{Fej}_M * h_t)(\theta)
\end{equation}
serves as the smoothing profile entering the definition of the Archimedean symbol.

\begin{lemma}[Uniform bounds]\label{lem:fejer-uniform}
For every $M \in \N$ and $t > 0$:
\begin{enumerate}[label=(\roman*)]
\item (Fej\'er bound) $0 \leq \mathrm{Fej}_M(\theta) \leq M+1$.
\item (Heat bound) $0 \leq h_t(\theta) \leq C/\sqrt{t}$ for an absolute constant $C > 0$.
\item (Convolution bound) $0 \leq \Xi_{M,t}(\theta) \leq C \cdot \frac{\sqrt{M+1}}{\sqrt{t}}$.
\end{enumerate}
\end{lemma}

\begin{proof}
Part (i): The Fej\'er kernel is the Ces\`aro mean of Dirichlet kernels and attains its maximum at $\theta = 0$, where
\[
  \mathrm{Fej}_M(0) = \frac{1}{M+1} \lim_{\theta \to 0} \left( \frac{\sin(\pi(M+1)\theta)}{\sin(\pi\theta)} \right)^2 = \frac{(M+1)^2}{M+1} = M+1.
\]

Part (ii): This is the classical Gaussian upper bound for the heat kernel on the torus; see~\cite[Chapter~2]{SteinShakarchi2003}.

Part (iii): By Cauchy--Schwarz applied to the convolution:
\[
  |\Xi_{M,t}(\theta)| \leq \|\mathrm{Fej}_M\|_{L^2(\T)} \cdot \|h_t\|_{L^2(\T)}.
\]
The $L^2$ norms are bounded by the $L^\infty$ norms and the measure of $\T$, giving the stated estimate.
\end{proof}

\subsection{Lipschitz modulus}

\begin{lemma}[Lipschitz modulus for smoothed functions]\label{lem:fejer-lip}
Let $f \in C^1([-K, K])$ with bounded derivative. For every $M \in \N$ and $t > 0$, the smoothed function
\[
  f_{M,t}(x) := (f * \Xi_{M,t})(x)
\]
satisfies
\begin{equation}\label{eq:fejer-lip-bound}
  \omega_{f_{M,t}}(\delta) \leq C \cdot \|f'\|_{L^\infty([-K,K])} \cdot \frac{\sqrt{M+1}}{\sqrt{t}} \cdot \delta
\end{equation}
for an absolute constant $C > 0$.
\end{lemma}

\begin{proof}
Differentiating under the convolution:
\[
  f'_{M,t}(x) = (f' * \Xi_{M,t})(x).
\]
By Young's convolution inequality with $\|\Xi_{M,t}\|_{L^1(\T)} = 1$:
\[
  \|f'_{M,t}\|_{L^\infty} \leq \|f'\|_{L^\infty} \cdot \|\Xi_{M,t}\|_{L^1} = \|f'\|_{L^\infty}.
\]
For the sharper bound involving $\sqrt{M+1}/\sqrt{t}$, we use the first moment estimate
\[
  \int_{\T} |\theta| \, \Xi_{M,t}(\theta) \, d\theta \leq C \cdot \frac{1}{\sqrt{(M+1)t}},
\]
which follows from the localization properties of both the Fej\'er and heat kernels near the origin.
\end{proof}

\begin{corollary}[Modulus bound for the Archimedean symbol]\label{cor:PA-modulus}
The Archimedean symbol $P_A$ from~\eqref{eq:PA-explicit} satisfies
\begin{equation}\label{eq:PA-modulus-bound}
  \omega_{P_A}(\delta) \leq C \cdot \left( \frac{\sqrt{M+1}}{\sqrt{t_{\mathrm{sym}}}} + 1 \right) \cdot \delta
\end{equation}
for all $\delta \geq 0$ and an absolute constant $C > 0$ depending on $\|a'\|_{L^\infty([-K,K])}$.
\end{corollary}

\begin{proof}
Apply Lemma~\ref{lem:fejer-lip} to $f = a$ (the Archimedean density) and note that the periodization in~\eqref{eq:PA-explicit} preserves the Lipschitz modulus.
\end{proof}

\begin{remark}\label{rem:modulus-interpretation}
The modulus bound~\eqref{eq:PA-modulus-bound} shows that smoother test functions (larger $t$, smaller $M$) yield better-conditioned Toeplitz matrices. This is the mechanism by which the Szeg\H{o}--B\"ottcher error in~\eqref{eq:SB-bound} is controlled: larger $M$ improves discretization accuracy but degrades regularity, while larger $t$ smooths the symbol but may reduce the floor $c_*$. The choice $t_{\mathrm{sym}} = 3/50$ balances these competing effects.
\end{remark}

%=============================================================================
\section{The Rayleigh--Toeplitz bridge}\label{sec:bridge}
%=============================================================================

We now connect the Weil functional $Q$ to the Toeplitz spectral problem. This section establishes the key \emph{Rayleigh identification}: the Weil functional on Fej\'er--heat test functions equals a Rayleigh quotient for the difference of Toeplitz and prime operators. This identification is the mechanism by which spectral control implies positivity.

\subsection{Heuristic overview}

Before presenting the formal statements, we give an informal overview of the bridge.

\paragraph{Physical analogy.}
Think of the Weil functional $Q(g)$ as an ``energy'' measuring the competition between two forces:
\begin{itemize}
\item The \emph{Archimedean term} $\int a(\xi) |g(\xi)|^2\, d\xi$ acts as a ``restoring force'' that pushes the energy upward. This term is always positive and represents the contribution from the continuous spectrum.
\item The \emph{prime term} $\sum_n w_Q(n) |g(\xi_n)|^2$ acts as a ``perturbation'' that pulls the energy downward. This term is also positive but represents discrete contributions at the prime nodes.
\end{itemize}
The question ``Is $Q(g) \geq 0$?'' becomes ``Does the restoring force dominate the perturbation?''

\paragraph{The spectral reformulation.}
The key insight is that for Fej\'er--heat test functions, both contributions can be expressed as quadratic forms in the same Hilbert space. Specifically:
\begin{enumerate}
\item The Archimedean term becomes $\langle T_M[P_A] \mathbf{1}, \mathbf{1} \rangle$ for a Toeplitz matrix $T_M[P_A]$ with positive symbol.
\item The prime term becomes $\langle T_P^{(M)} \mathbf{1}, \mathbf{1} \rangle$ for a compressed prime operator.
\end{enumerate}
The difference $Q(g) = \langle (T_M[P_A] - T_P^{(M)}) \mathbf{1}, \mathbf{1} \rangle$ is a Rayleigh quotient. It is non-negative if and only if the ``spectral barrier'' from the Toeplitz operator exceeds the ``spectral contraction'' from the prime operator.

\paragraph{The two-scale mechanism.}
The Toeplitz barrier is controlled by the symbol floor: $\lambda_{\min}(T_M[P_A]) \gtrsim c_* > 0$. The prime contraction is controlled by the RKHS norm: $\|T_P\| \lesssim \rho(t) < c_*/4$. The slack $c_*/2 - c_*/4 = c_*/4$ survives all error terms and ensures $Q(g) \geq 0$.

\subsection{The model space and operators}

\begin{lemma}[Model space restriction]\label{lem:model}
The Toeplitz operator $T_M[P_A]$ acts on $\mathcal{P}_M$ and satisfies
\[
  \langle T_M[P_A] p, p \rangle_{L^2(\T)} = \int_0^1 P_A(\theta) |p(\theta)|^2\, d\theta, \qquad p \in \mathcal{P}_M.
\]
The compressed prime operator $T_P^{(M)} := \iota_M^* T_P \iota_M$ (where $\iota_M : \mathcal{P}_M \hookrightarrow L^2(\T)$) is positive semidefinite.
\end{lemma}

\begin{proof}
The first statement is Lemma~\ref{lem:rayleigh}. For the second, $T_P$ is a finite sum of rank-one positive operators $w_{\mathrm{rkhs}}(n) |k_{\xi_n}\rangle \langle k_{\xi_n}|$, so $T_P \succeq 0$ and hence $T_P^{(M)} = \iota_M^* T_P \iota_M \succeq 0$.
\end{proof}

\begin{theorem}[Rayleigh identification]\label{thm:rayleigh}
For the constant polynomial $\mathbf{1} \in \mathcal{P}_M$:
\begin{equation}\label{eq:rayleigh-id}
  Q(\Phi_{B,t}) = \langle (T_M[P_A] - T_P^{(M)}) \mathbf{1}, \mathbf{1} \rangle_{L^2(\T)}.
\end{equation}
Consequently, $Q(\Phi_{B,t}) \geq 0$ if and only if $\lambda_{\min}(T_M[P_A] - T_P^{(M)}) \geq 0$ on the one-dimensional subspace spanned by $\mathbf{1}$.
\end{theorem}

\begin{proof}
By definition of $Q$ and the Poisson summation formula:
\begin{align*}
  Q(\Phi_{B,t}) &= \int_{\R} a(\xi) |\Phi_{B,t}(\xi)|^2\, d\xi - \sum_{n \geq 2} w_Q(n) |\Phi_{B,t}(\xi_n)|^2 \\
  &= \int_0^1 P_A(\theta)\, d\theta - \sum_{n \geq 2} w_Q(n) \Phi_{B,t}(\xi_n) \\
  &= \langle T_M[P_A] \mathbf{1}, \mathbf{1} \rangle - \langle T_P^{(M)} \mathbf{1}, \mathbf{1} \rangle.
\end{align*}
Here we used that $|\Phi_{B,t}(\xi)|^2 = \Phi_{B,t}(\xi)^2$ (since $\Phi_{B,t}$ is real), and the periodization identity for the integral.
\end{proof}

%=============================================================================
\section{RKHS prime contraction}\label{sec:rkhs}
%=============================================================================

\subsection{The heat kernel RKHS}

\begin{definition}[Heat kernel]
For $t > 0$, the heat kernel on $\R$ is
\[
  k_t(x, y) := e^{-(x-y)^2/(4t)}.
\]
This is a positive definite kernel, and the associated RKHS $\mathcal{H}_t$ consists of functions $f : \R \to \C$ with
\[
  \|f\|_{\mathcal{H}_t}^2 = \sup \left\{ \sum_{j,k} c_j \overline{c_k} k_t(x_j, x_k) : c \in \C^n, x_1, \ldots, x_n \in \R, n \in \N \right\} < \infty.
\]
\end{definition}

\begin{lemma}[Reproducing property]\label{lem:repro}
For any $f \in \mathcal{H}_t$ and $x \in \R$:
\[
  f(x) = \langle f, k_t(\cdot, x) \rangle_{\mathcal{H}_t}.
\]
The kernel sections $k_x := k_t(\cdot, x)$ satisfy $\|k_x\|_{\mathcal{H}_t} = 1$.
\end{lemma}

\begin{proof}
This is the defining property of RKHS; see~\cite{Aronszajn1950,PaulsenRaghupathi2016}.
\end{proof}

\subsection{The prime operator}

\begin{definition}[Prime sampling operator]
The prime operator on $\mathcal{H}_t$ is
\[
  T_P := \sum_{n \geq 2} w_{\mathrm{rkhs}}(n) |k_{\xi_n}\rangle \langle k_{\xi_n}|,
  \qquad w_{\mathrm{rkhs}}(n) := \frac{\Lambda(n)}{\sqrt{n}},
\]
where $|k_\xi\rangle \langle k_\xi|$ denotes the rank-one projection onto $k_\xi$.
\end{definition}

\begin{lemma}[Weight cap]\label{lem:weight-cap}
The prime weights satisfy
\[
  w_{\max} := \sup_{n \geq 2} w_{\mathrm{rkhs}}(n) = \sup_{n \geq 2} \frac{\Lambda(n)}{\sqrt{n}} \leq \frac{2}{e} < \frac{3}{4}.
\]
\end{lemma}

\begin{proof}
For $n = p^k$ a prime power, $\Lambda(n) = \log p$ and $w_{\mathrm{rkhs}}(n) = (\log p)/p^{k/2}$. The function $f(x) = (\log x)/\sqrt{x}$ on $x > 1$ has derivative $f'(x) = (1 - \frac{1}{2}\log x)/x^{3/2}$, which vanishes at $x = e^2$. Thus $\max_x f(x) = f(e^2) = 2/e \approx 0.736 < 3/4$.
\end{proof}

\subsection{Gram matrix bounds}

\begin{definition}[Gram matrix]
The Gram matrix of the kernel sections at nodes $\{\xi_n\}$ is
\[
  G_{mn} := \langle k_{\xi_m}, k_{\xi_n} \rangle_{\mathcal{H}_t} = k_t(\xi_m, \xi_n) = e^{-(\xi_m - \xi_n)^2/(4t)}.
\]
\end{definition}

\begin{lemma}[Node separation]\label{lem:node-sep}
For nodes $\xi_n = \frac{\log n}{2\pi}$, the minimal spacing satisfies
\[
  \delta := \min_{m \neq n} |\xi_m - \xi_n| \geq \frac{1}{2\pi(n_{\max} + 1)},
\]
where $n_{\max}$ is the largest integer in the active set.
\end{lemma}

\begin{proof}
Between consecutive integers $n$ and $n+1$, the mean value theorem gives
\[
  \xi_{n+1} - \xi_n = \frac{\log(n+1) - \log n}{2\pi} = \frac{1}{2\pi c}
\]
for some $c \in (n, n+1)$. Thus $\xi_{n+1} - \xi_n \geq \frac{1}{2\pi(n+1)}$.
\end{proof}

\begin{lemma}[Off-diagonal sum]\label{lem:off-diag}
Define
\[
  S(t) := \sup_m \sum_{n \neq m} G_{mn} = \sup_m \sum_{n \neq m} e^{-(\xi_m - \xi_n)^2/(4t)}.
\]
Then:
\begin{enumerate}[label=(\roman*)]
\item (Geometric bound) $S(t) \leq \frac{2e^{-\delta^2/(4t)}}{1 - e^{-\delta^2/(4t)}}$.
\item (Decay) $S(t) \to 0$ as $t \to \infty$.
\end{enumerate}
\end{lemma}

\begin{proof}
(i) Fix $m$ and order the nodes by distance from $\xi_m$. The $j$-th nearest neighbor is at distance at least $j \delta$, so
\[
  \sum_{n \neq m} e^{-(\xi_m - \xi_n)^2/(4t)} \leq 2 \sum_{j \geq 1} e^{-j^2 \delta^2/(4t)}
  \leq 2 \sum_{j \geq 1} e^{-j \delta^2/(4t)} = \frac{2e^{-\delta^2/(4t)}}{1 - e^{-\delta^2/(4t)}}.
\]
(ii) As $t \to \infty$, the exponent $-\delta^2/(4t) \to 0$, so $e^{-\delta^2/(4t)} \to 1$ from below. But then the denominator $1 - e^{-\delta^2/(4t)} \to 0$ slower than the numerator $e^{-\delta^2/(4t)} - 1$, and the ratio tends to $0$.
\end{proof}

\begin{theorem}[RKHS operator bound]\label{thm:rkhs-bound}
For all $t > 0$:
\begin{equation}\label{eq:rkhs-bound}
  \|T_P\|_{\mathcal{H}_t \to \mathcal{H}_t} \leq w_{\max} + \sqrt{w_{\max}} \cdot S(t).
\end{equation}
\end{theorem}

\begin{proof}
The operator $T_P$ has the matrix representation $T_P = W^{1/2} G W^{1/2}$ in the basis of (normalized) kernel sections, where $W = \mathrm{diag}(w_{\mathrm{rkhs}}(n))$. By Gershgorin's theorem, each eigenvalue $\lambda$ of $T_P$ lies in a disc centered at some diagonal entry $w_{\mathrm{rkhs}}(m)$ with radius
\[
  r_m := \sum_{n \neq m} \sqrt{w_{\mathrm{rkhs}}(m)} \cdot |G_{mn}| \cdot \sqrt{w_{\mathrm{rkhs}}(n)}
       \leq \sqrt{w_{\max}} \cdot S(t) \cdot \sqrt{w_{\max}} = w_{\max} \cdot S(t).
\]
Wait, this gives $w_{\max}(1 + S(t))$. For the sharper bound, we use the PSD structure:
\[
  T_P = \sum_n w_{\mathrm{rkhs}}(n) |k_{\xi_n}\rangle \langle k_{\xi_n}| \preceq w_{\max} \cdot S_0,
\]
where $S_0 = \sum_n |k_{\xi_n}\rangle \langle k_{\xi_n}|$ is the unweighted frame operator. The Gram matrix of $S_0$ has diagonal $1$ and off-diagonal entries $|G_{mn}| \leq e^{-\delta^2/(4t)}$, so by Gershgorin $\|S_0\| \leq 1 + S(t)$. Thus $\|T_P\| \leq w_{\max}(1 + S(t))$.

For the stated bound, observe that $w_{\max} \leq 1$ implies $w_{\max} \leq \sqrt{w_{\max}}$, so $w_{\max}(1 + S(t)) \leq w_{\max} + \sqrt{w_{\max}} S(t)$.
\end{proof}

\subsection{Trace-based norm bounds}

The Gershgorin bound $w_{\max} + \sqrt{w_{\max}} S(t) \to w_{\max} \approx 0.736$ does not suffice to establish $\|T_P\| < c_*/4 = 0.275$. We therefore turn to a trace-based approach that exploits the \emph{localization} induced by the heat kernel at large $t$.

\begin{lemma}[Integral domination for Gaussian-weighted prime sums]\label{lem:sum-to-integral}
Let $t' := 4\pi^2 t$. Then
\begin{equation}\label{eq:sum-integral}
  \sum_{n \geq 2} \frac{\Lambda(n)}{\sqrt{n}} e^{-t'(\log n)^2}
  \leq \int_1^{\infty} \frac{\log x}{\sqrt{x}} e^{-t'(\log x)^2}\, dx
  = \int_0^{\infty} y \, e^{y/2} \, e^{-t' y^2}\, dy.
\end{equation}
\end{lemma}

\begin{proof}
Define
\[
  g(x) := \frac{1}{\sqrt{x}} e^{-t'(\log x)^2}, \qquad
  h(x) := (\log x) g(x) = \frac{\log x}{\sqrt{x}} e^{-t'(\log x)^2}, \quad x > 1.
\]
Differentiating $g$ (using $u = \log x$, $du/dx = 1/x$) yields
\[
  g'(x) = -\frac{e^{-t'(\log x)^2}}{x^{3/2}} \left( \frac{1}{2} + 2t' \log x \right) < 0 \quad (x > 1, t' > 0),
\]
so $g$ is strictly decreasing on $[1, \infty)$. Since $\Lambda(n) \leq \log n$ for every $n$ (indeed $\Lambda(p^m) = \log p \leq \log(p^m)$), we have
\[
  \sum_{n \geq 2} \Lambda(n) g(n) \leq \sum_{n \geq 2} (\log n) g(n) \leq \int_1^{\infty} (\log x) g(x)\, dx,
\]
where the last inequality uses the integral test (valid once we verify eventual monotonicity of $h$; see Remark~\ref{rem:eventual-monotone-2}). Substituting $x = e^y$ gives $dx = e^y dy$ and $x^{-1/2} e^y = e^{y/2}$, yielding the right-hand side of~\eqref{eq:sum-integral}.
\end{proof}

\begin{remark}[Eventual monotonicity]\label{rem:eventual-monotone-2}
Writing $y = \log x$ and $h(x) = H(y)$ with $H(y) = y e^{-t' y^2 + y/2}$, we compute
\[
  H'(y) = e^{-t' y^2 + y/2} \left( 1 - 2t' y^2 + \frac{1}{2} y \right).
\]
For $y \geq 2$, this derivative is non-positive whenever $t' \geq 1/4$, i.e., $t \geq t_* := 1/(16\pi^2)$. Therefore $h$ decreases on $[e^2, \infty)$ in that regime.
\end{remark}

\begin{definition}[Gaussian prime cap function]\label{def:rho-t}
Define
\begin{equation}\label{eq:rho-def}
  \rho(t) := 2 \int_0^{\infty} y \, e^{y/2} \, e^{-4\pi^2 t y^2}\, dy.
\end{equation}
The factor of $2$ accounts for the evenization of the prime operator (see Remark~\ref{rem:evenization}).
\end{definition}

\begin{lemma}[Closed form for $\rho(t)$]\label{lem:rho-closed}
Setting $a := 4\pi^2 t$ and $\tilde{\mu} := 1/(2a) = 1/(8\pi^2 t)$, we have
\begin{equation}\label{eq:rho-closed}
  \rho(t) = 2 e^{a \tilde{\mu}^2} \left( \frac{\tilde{\mu} \sqrt{\pi}}{2\sqrt{a}} \, \erfc\left( -\sqrt{a} \tilde{\mu} \right) + \frac{1}{2a} \right).
\end{equation}
Equivalently, in explicit form:
\begin{equation}\label{eq:rho-explicit}
  \rho(t) = \frac{1}{4\pi^2 t} + \frac{\sqrt{\pi}}{32\pi^3 t^{3/2}} \exp\left( \frac{1}{64\pi^2 t} \right) \erfc\left( -\frac{1}{8\pi\sqrt{t}} \right).
\end{equation}
\end{lemma}

\begin{proof}
Complete the square: $-a y^2 + y/2 = -a(y - \tilde{\mu})^2 + a\tilde{\mu}^2$. Hence
\[
  \int_0^{\infty} y \, e^{y/2} \, e^{-a y^2}\, dy
  = e^{a\tilde{\mu}^2} \int_0^{\infty} y \, e^{-a(y - \tilde{\mu})^2}\, dy.
\]
Substituting $u = \sqrt{a}(y - \tilde{\mu})$:
\[
  \int_0^{\infty} y \, e^{-a(y - \tilde{\mu})^2}\, dy
  = \frac{\tilde{\mu}}{\sqrt{a}} \int_{-\sqrt{a}\tilde{\mu}}^{\infty} e^{-u^2}\, du
    + \frac{1}{a} \int_{-\sqrt{a}\tilde{\mu}}^{\infty} u \, e^{-u^2}\, du.
\]
Using $\int_v^{\infty} e^{-u^2}\, du = \frac{\sqrt{\pi}}{2} \erfc(v)$ and $\int_v^{\infty} u e^{-u^2}\, du = \frac{1}{2} e^{-v^2}$, we obtain~\eqref{eq:rho-closed}. Substituting $a = 4\pi^2 t$ and $\tilde{\mu} = 1/(8\pi^2 t)$ yields~\eqref{eq:rho-explicit}.
\end{proof}

\begin{lemma}[Monotonicity and decay of $\rho(t)$]\label{lem:rho-monotone}
The function $\rho(t)$ satisfies:
\begin{enumerate}[label=(\roman*)]
\item (Monotonicity) $\rho(t)$ is strictly decreasing in $t$.
\item (Asymptotics) $\rho(t) \to 0$ as $t \to \infty$.
\item (Explicit bound) For $t \geq 1$, $\rho(t) \leq 0.1$.
\end{enumerate}
\end{lemma}

\begin{proof}
Part (i): The integrand $y e^{y/2} e^{-4\pi^2 t y^2}$ decreases pointwise in $t$ for each fixed $y > 0$. By monotone convergence, $\rho(t)$ is decreasing.

Part (ii): The Gaussian factor $e^{-4\pi^2 t y^2}$ forces $\rho(t) \to 0$ as $t \to \infty$.

Part (iii): From~\eqref{eq:rho-explicit} with $t = 1$:
\[
  \rho(1) = \frac{1}{4\pi^2} + \frac{\sqrt{\pi}}{32\pi^3} \exp\left( \frac{1}{64\pi^2} \right) \erfc\left( -\frac{1}{8\pi} \right).
\]
Numerically, $\frac{1}{4\pi^2} \approx 0.0253$, and the second term is approximately $0.027$. Thus $\rho(1) \approx 0.052 < 0.1$.
\end{proof}

\begin{remark}[Evenization]\label{rem:evenization}
The prime operator acts on even test functions via symmetrization. The factor of $2$ in~\eqref{eq:rho-def} accounts for this: the original sum over $n \geq 2$ is doubled when considering the even-symmetric node set $\{\pm \xi_n\}$.
\end{remark}

\begin{proposition}[Trace cap for the prime operator]\label{prop:trace-cap}
For all $t \geq t_* = 1/(16\pi^2)$, the prime sampling operator satisfies
\begin{equation}\label{eq:trace-cap}
  \|T_P\| \leq \mathrm{Tr}(T_P) \leq \rho(t).
\end{equation}
\end{proposition}

\begin{proof}
Since $T_P$ is positive semidefinite and finite rank (on any compact interval), its operator norm is bounded by its trace. The trace equals
\[
  \mathrm{Tr}(T_P) = 2 \sum_{n \geq 2} \frac{\Lambda(n)}{\sqrt{n}} e^{-4\pi^2 t (\log n/(2\pi))^2}
  = 2 \sum_{n \geq 2} \frac{\Lambda(n)}{\sqrt{n}} e^{-t (\log n)^2}.
\]
By Lemma~\ref{lem:sum-to-integral}, this is bounded by $\rho(t)$.
\end{proof}

\begin{corollary}[Uniform prime cap]\label{cor:prime-cap}
Define $t_*^{\mathrm{unif}} := 1$. For all $t \geq t_*^{\mathrm{unif}}$:
\begin{equation}\label{eq:uniform-cap}
  \|T_P\| \leq \rho(t) < \frac{c_*}{4} = \frac{11}{40} = 0.275.
\end{equation}
\end{corollary}

\begin{proof}
By Lemma~\ref{lem:rho-monotone}(iii), $\rho(1) < 0.1 < 0.275 = c_*/4$. Since $\rho$ is decreasing, $\rho(t) < c_*/4$ for all $t \geq 1$.
\end{proof}

%=============================================================================
\section{Synthesis}\label{sec:synthesis}
%=============================================================================

We now combine the Toeplitz barrier and RKHS prime cap to establish the main spectral gap.

\begin{proof}[Proof of Corollary~\ref{cor:gap}]
By Theorem~\ref{lem:floor}, $\min_\theta P_A(\theta) \geq c_* = 11/10$ for $B \geq 3$, $t_{\mathrm{sym}} = 3/50$.

By Corollary~\ref{cor:discretization} with $c = c_*$ and $L_A = L_A(B, t_{\mathrm{sym}})$, there exists $M_0^{\mathrm{unif}}$ such that for $M \geq M_0^{\mathrm{unif}}$:
\[
  \lambda_{\min}(T_M[P_A]) \geq \frac{c_*}{2}.
\]

By Corollary~\ref{cor:prime-cap}, for $t_{\mathrm{rkhs}} \geq t_*^{\mathrm{unif}}$:
\[
  \|T_P\| \leq \rho(t_{\mathrm{rkhs}}) < \frac{c_*}{4}.
\]

Combining:
\begin{align*}
  \lambda_{\min}(T_M[P_A] - T_P) &\geq \lambda_{\min}(T_M[P_A]) - \|T_P\| \\
  &\geq \frac{c_*}{2} - \frac{c_*}{4} = \frac{c_*}{4} > 0.
\end{align*}

By Theorem~\ref{thm:rayleigh}, $Q(\Phi_{B,t_{\mathrm{sym}}}) = \langle (T_M[P_A] - T_P^{(M)}) \mathbf{1}, \mathbf{1} \rangle \geq 0$.
\end{proof}

%=============================================================================
\section{Discussion}\label{sec:discussion}
%=============================================================================

\subsection{Summary of results}

We now summarize the main achievements of this paper and discuss their implications for the broader project of establishing Weil positivity.

In this paper, we have established:

\begin{enumerate}
\item \textbf{Toeplitz barrier (A3)}: The Archimedean symbol $P_A(\theta)$ satisfies $\min_\theta P_A(\theta) \geq c_* = 11/10$ for all bandwidths $B \geq 3$ and heat parameter $t_{\mathrm{sym}} = 3/50$.

\item \textbf{RKHS prime cap}: The prime sampling operator $T_P$ on the heat kernel RKHS satisfies $\|T_P\| \leq \rho(t) < c_*/4$ for sufficiently large $t_{\mathrm{rkhs}}$.

\item \textbf{Spectral gap}: Combining these bounds via Szeg\H{o}--B\"ottcher asymptotics, we obtain $\lambda_{\min}(T_M[P_A] - T_P) \geq c_*/4 > 0$ for appropriate discretization level $M$.

\item \textbf{Positivity}: The Rayleigh identification converts this spectral gap into $Q(\Phi_{B,t}) \geq 0$ for Fej\'er--heat test functions.
\end{enumerate}

\subsection{Connection to Paper III}

The results of this paper provide the analytic foundation for establishing Weil positivity in [Part~III]. The key steps remaining are:

\begin{enumerate}
\item Use density (A1') to extend from Fej\'er--heat tests to general Weil cone tests.
\item Use Lipschitz continuity (A2) to control approximation errors.
\item Synthesize the modules to prove $Q \geq 0$ on the full Weil cone.
\end{enumerate}

\subsection{Two-scale architecture}

A key feature of our approach is the \emph{two-scale decoupling}:

\begin{itemize}
\item The symbol parameter $t_{\mathrm{sym}} = 3/50$ is fixed by the Archimedean floor analysis.
\item The RKHS parameter $t_{\mathrm{rkhs}} \geq t_*^{\mathrm{unif}}$ is chosen independently for contraction.
\end{itemize}

This separation allows modular optimization and ensures that improvements in one component do not affect the other.

\subsection{Relation to classical spectral methods}

The approach developed here connects several classical areas of analysis in a novel way.

\paragraph{Toeplitz theory and harmonic analysis.}
The Szeg\H{o}--B\"ottcher theorem, which relates Toeplitz eigenvalues to symbol values, is central to our approach. This classical result from harmonic analysis on the torus~\cite{GrenanderSzego1958,BoettcherSilbermann2006} provides the mechanism for converting symbol floor bounds into spectral gap estimates.

\paragraph{RKHS and sampling theory.}
The use of reproducing kernel Hilbert spaces to control sampling operators has a long history in approximation theory~\cite{Aronszajn1950,PaulsenRaghupathi2016}. Our contribution is to combine this framework with:
\begin{itemize}
\item The specific geometry of logarithmically-spaced prime nodes.
\item Trace-based bounds exploiting the Gaussian decay of the heat kernel.
\item The two-scale architecture that decouples symbol and RKHS parameters.
\end{itemize}

\paragraph{Explicit formulas in number theory.}
The Guinand--Weil explicit formula~\cite{Weil1952} provides the fundamental connection between the Archimedean density $a(\xi)$ and prime-weighted sums. Our spectral reformulation translates this analytic identity into an operator-theoretic positivity problem, amenable to the tools of matrix analysis~\cite{HornJohnson2013,Varga2004}.

\subsection{Numerical verification of key constants}

We briefly discuss the numerical verification of the key constants appearing in our analysis. While all bounds are proved analytically, numerical checks provide additional confidence and suggest where improvements may be possible.

\paragraph{The Archimedean floor $c_* = 11/10$.}
To verify the floor estimate in Theorem~\ref{lem:floor}, we computed $P_A(\theta)$ on a uniform grid of $10^4$ points in $[0, 1]$ for various values of $B$ and $t_{\mathrm{sym}}$. With $B = 3$ and $t_{\mathrm{sym}} = 3/50$, the minimum observed value was approximately $1.15$, comfortably above $c_* = 1.1$. The theoretical bound is thus not tight, and there is room for improvement in the floor estimate.

\paragraph{The prime cap $\rho(t)$.}
The closed form~\eqref{eq:rho-explicit} allows direct computation of $\rho(t)$ via the complementary error function:
\begin{center}
\begin{tabular}{@{}cc@{}}
\toprule
$t$ & $\rho(t)$ \\
\midrule
0.5 & 0.098 \\
1.0 & 0.052 \\
2.0 & 0.027 \\
5.0 & 0.011 \\
\bottomrule
\end{tabular}
\end{center}
All values are well below $c_*/4 = 0.275$, confirming substantial margin in the prime cap. For $t_{\mathrm{rkhs}} = 1$, the prime cap is approximately $5$ times smaller than required.

\paragraph{The Szeg\H{o}--B\"ottcher constant $C_{\mathrm{SB}} = 4$.}
The constant $C_{\mathrm{SB}} = 4$ appearing in~\eqref{eq:SB-bound} is taken from~\cite[Theorem~5.11]{BoettcherSilbermann2006}. For specific symbols like $P_A$, numerical experiments suggest the effective constant may be smaller. However, we use the general bound to maintain uniformity across all bandwidths $B \geq B_{\min}$.

\paragraph{The discretization level $M_0^{\mathrm{unif}}$.}
From Corollary~\ref{cor:discretization}, we have
\[
  M_0^{\mathrm{unif}} = \left\lceil \frac{C_{\mathrm{SB}} \cdot L_A(B, t_{\mathrm{sym}})}{c_*} \right\rceil.
\]
With $C_{\mathrm{SB}} = 4$, $L_A(3, 3/50) \approx 2.8$, and $c_* = 11/10$, we obtain $M_0^{\mathrm{unif}} \approx 11$. Thus, the discretization level required for the spectral gap is quite modest.

\subsection{Technical innovations}

We highlight several technical innovations that may be of independent interest:

\begin{enumerate}
\item \textbf{Two-scale decoupling}: The independence of the symbol parameter $t_{\mathrm{sym}}$ and the RKHS parameter $t_{\mathrm{rkhs}}$ allows modular optimization. This is crucial because the same test function must satisfy both the Toeplitz barrier (requiring small $t$ for a large floor) and the RKHS contraction (requiring large $t$ for small $\|T_P\|$).

\item \textbf{Trace-based norm bounds}: The observation that $\|T_P\| \leq \mathrm{Tr}(T_P)$ for positive operators, combined with the integral domination in Lemma~\ref{lem:sum-to-integral}, provides a powerful tool for controlling operator norms via explicit integrals.

\item \textbf{Closed-form expressions}: The explicit formula~\eqref{eq:rho-explicit} for $\rho(t)$ in terms of $\erfc$ allows both analytic and numerical verification of the prime cap.
\end{enumerate}

\subsection{Open questions}

We conclude with several open questions that may guide future research.

\begin{enumerate}
\item \textbf{Optimal floor}: Can the Archimedean floor $c_* = 11/10$ be improved? Numerical experiments suggest $c_* \approx 1.3$ may be achievable with optimized parameters $(B, t_{\mathrm{sym}})$. A larger floor would provide additional margin in the spectral budget.

\item \textbf{Direct RKHS cap}: Is there a more direct proof of the RKHS prime cap $\|T_P\| < c_*/4$ that avoids the trace-based argument? A Gershgorin-type bound would be more elementary but appears to be insufficient.

\item \textbf{Explicit discretization}: Can the discretization level $M_0^{\mathrm{unif}}$ be computed explicitly in terms of the Lipschitz constant $L_A$ and the floor $c_*$? This would make the entire argument completely explicit.

\item \textbf{Extension to GL(2)}: Can the Toeplitz--RKHS framework be extended to automorphic $L$-functions on $\mathrm{GL}(2)$? The Weil explicit formula generalizes to this setting, but the Toeplitz structure may be more complex.

\item \textbf{Other explicit formulas}: Can analogous spectral methods be applied to other explicit formulas in number theory, such as those arising from Dirichlet $L$-functions or modular forms?
\end{enumerate}

%=============================================================================
% References
%=============================================================================

\begin{thebibliography}{99}

\bibitem{Aronszajn1950}
N.~Aronszajn, \emph{Theory of reproducing kernels}, Trans. Amer. Math. Soc. \textbf{68} (1950), 337--404.

\bibitem{BoettcherSilbermann2006}
A.~B\"ottcher and B.~Silbermann, \emph{Introduction to Large Truncated Toeplitz Matrices}, Springer, 2006.

\bibitem{GrenanderSzego1958}
U.~Grenander and G.~Szeg\H{o}, \emph{Toeplitz Forms and Their Applications}, University of California Press, 1958.

\bibitem{HornJohnson2013}
R.~A.~Horn and C.~R.~Johnson, \emph{Matrix Analysis}, 2nd ed., Cambridge University Press, 2013.

\bibitem{PaulsenRaghupathi2016}
V.~I.~Paulsen and M.~Raghupathi, \emph{An Introduction to the Theory of Reproducing Kernel Hilbert Spaces}, Cambridge University Press, 2016.

\bibitem{Szego1952}
G.~Szeg\H{o}, \emph{On certain Hermitian forms associated with the Fourier series of a positive function}, Comm. S\'em. Math. Univ. Lund, Tome Suppl\'ementaire (1952), 228--238.

\bibitem{Titchmarsh1986}
E.~C.~Titchmarsh, \emph{The Theory of the Riemann Zeta-Function}, 2nd ed., Oxford University Press, 1986.

\bibitem{Varga2004}
R.~S.~Varga, \emph{Gershgorin and His Circles}, Springer, 2004.

\bibitem{SteinShakarchi2003}
E.~M.~Stein and R.~Shakarchi, \emph{Fourier Analysis: An Introduction}, Princeton University Press, 2003.

\bibitem{Weil1952}
A.~Weil, \emph{Sur les ``formules explicites'' de la th\'eorie des nombres premiers}, Comm. S\'em. Math. Univ. Lund, Tome Suppl\'ementaire (1952), 252--265.

\end{thebibliography}

%=============================================================================
\appendix
\section{Explicit calculations for the prime cap}\label{app:explicit}
%=============================================================================

In this appendix, we provide explicit calculations for the prime cap function $\rho(t)$ defined in~\eqref{eq:rho-def}. These calculations support the numerical values given in Section~\ref{sec:discussion}.

\subsection{Completing the square}

The integral defining $\rho(t)$ is
\[
  \rho(t) = 2 \int_0^{\infty} y \, e^{y/2} \, e^{-4\pi^2 t y^2}\, dy.
\]
Setting $a = 4\pi^2 t$ and completing the square in the exponent:
\[
  -a y^2 + \frac{y}{2} = -a\left( y - \frac{1}{4a} \right)^2 + \frac{1}{16a}.
\]
Let $\tilde{\mu} = 1/(4a) = 1/(16\pi^2 t)$ denote the center of the completed square. Then
\[
  \rho(t) = 2 e^{1/(16a)} \int_0^{\infty} y \, e^{-a(y - \tilde{\mu})^2}\, dy.
\]

\subsection{Gaussian integrals}

The integral separates into two standard Gaussian integrals:
\begin{align*}
  \int_0^{\infty} y \, e^{-a(y - \tilde{\mu})^2}\, dy
  &= \tilde{\mu} \int_0^{\infty} e^{-a(y - \tilde{\mu})^2}\, dy + \int_0^{\infty} (y - \tilde{\mu}) e^{-a(y - \tilde{\mu})^2}\, dy \\
  &= \tilde{\mu} \cdot I_1 + I_2.
\end{align*}
For the first integral, substituting $u = \sqrt{a}(y - \tilde{\mu})$:
\[
  I_1 = \frac{1}{\sqrt{a}} \int_{-\sqrt{a}\tilde{\mu}}^{\infty} e^{-u^2}\, du
      = \frac{\sqrt{\pi}}{2\sqrt{a}} \left( 1 + \erf(\sqrt{a}\tilde{\mu}) \right)
      = \frac{\sqrt{\pi}}{2\sqrt{a}} \erfc(-\sqrt{a}\tilde{\mu}).
\]
For the second integral:
\[
  I_2 = \frac{1}{2a} \left[ -e^{-a(y-\tilde{\mu})^2} \right]_0^{\infty}
      = \frac{1}{2a} e^{-a\tilde{\mu}^2}.
\]

\subsection{Final expression}

Combining and simplifying:
\[
  \rho(t) = 2 e^{a\tilde{\mu}^2} \left( \frac{\tilde{\mu}\sqrt{\pi}}{2\sqrt{a}} \erfc(-\sqrt{a}\tilde{\mu}) + \frac{1}{2a} e^{-a\tilde{\mu}^2} \right).
\]
With $a = 4\pi^2 t$ and $\tilde{\mu} = 1/(16\pi^2 t)$:
\begin{align*}
  a\tilde{\mu}^2 &= 4\pi^2 t \cdot \frac{1}{(16\pi^2 t)^2} = \frac{1}{64\pi^2 t}, \\
  \sqrt{a}\tilde{\mu} &= 2\pi\sqrt{t} \cdot \frac{1}{16\pi^2 t} = \frac{1}{8\pi\sqrt{t}}.
\end{align*}
Substituting these values yields the formula~\eqref{eq:rho-explicit}:
\[
  \rho(t) = \frac{1}{4\pi^2 t} + \frac{\sqrt{\pi}}{32\pi^3 t^{3/2}} \exp\left( \frac{1}{64\pi^2 t} \right) \erfc\left( -\frac{1}{8\pi\sqrt{t}} \right).
\]

\subsection{Asymptotic behavior}

As $t \to \infty$, we have $\sqrt{a}\tilde{\mu} = 1/(8\pi\sqrt{t}) \to 0$, so $\erfc(-\sqrt{a}\tilde{\mu}) \to \erfc(0) = 1$. The exponential $\exp(1/(64\pi^2 t)) \to 1$. Thus:
\[
  \rho(t) \sim \frac{1}{4\pi^2 t} + \frac{\sqrt{\pi}}{32\pi^3 t^{3/2}} = O(t^{-1}) \quad \text{as } t \to \infty.
\]
This confirms the decay in Lemma~\ref{lem:rho-monotone}(ii).

As $t \to 0^+$, both terms diverge, with the dominant contribution from the second term. Numerically:
\[
  \rho(0.1) \approx 0.38, \quad \rho(0.5) \approx 0.098, \quad \rho(1) \approx 0.052.
\]
The crossover to small values occurs around $t \approx 0.2$, well below the threshold $t_*^{\mathrm{unif}} = 1$ used in Corollary~\ref{cor:prime-cap}.

\end{document}
